% !TeX root = ../../book.tex
\section{集合运算}
\secbegin{secSetOperations}

在 \Cref{exPositiveNegativeSetBuilderNotation} 中，我们注意到 $[0,\infty)$ 是所有非负实数的集合。如果我们想讨论所有非负有理数的集合怎么办？如果能用 $[0,\infty)$ 和 $\mathbb{Q}$ 来表示这个集合就好了。

这就是\textit{集合运算}的用武之地 --- 它们允许我们使用先前定义的集合来引入新的集合。

\subsubsection*{交集 ($\cap$)}
两个集合的\textit{交集}是两个集合共同元素的集合。

\begin{definition}
\label{defIntersection}
\index{intersection!pairwise}
\nindex{intersection}{$\cap$}{intersection}
设 $X$ 和 $Y$ 为集合。 $X$ 和 $Y$ 的（\textbf{成对}）\textbf{交集}，表示为 $X \cap Y$ \inlatex{cap}\lindexmmc{cap}{$\cap$}，定义为
\[ X \cap Y = \{ a \mid a \in X \wedge a \in Y \} \]
\end{definition}

\begin{example}
根据交集的定义，我们有 $x \in [0,\infty) \cap \mathbb{Q}$ 当且仅当 $x \in [0,\infty)$ 且 $x \in \mathbb{Q} $。由于 $x \in [0,\infty)$ 当且仅当 $x$ 是非负实数（参见 \Cref{exPositiveNegativeSetBuilderNotation}），因此 $[0,\infty) \cap \mathbb{ Q}$ 是所有非负有理数的集合。
\end{example}

\begin{exercise}
证明 $[0,\infty) \cap \mathbb{Z} = \mathbb{N}$。
\end{exercise}

\begin{exercise}
写出下面集合的元素
\[ \{ 0, 1, 4, 7 \} \cap \{ 1, 2, 3, 4, 5 \} \]
\end{exercise}

\begin{exercise}
将 $[-2,5) \cap [4,7)$ 表示为单个区间。
\end{exercise}

\begin{proposition}
\label{propSubsetFromIntersection}
设 $X$ 和 $Y$ 为集合。证明 $X \subseteq Y$ 当且仅当 $X \cap Y = X$。
\end{proposition}

\begin{cproof}
%% BEGIN QUOTATION (xtrSoExample) %%
假设 $X \subseteq Y$。我们通过双重包含证明 $X \cap Y = X$。
\begin{itemize}
\item ($\subseteq$) 假设 $a \in X \cap Y$。则根据交集的定义 $a \in X$ 且 $a \in Y$，所以我们有 $a \in X$。
\item ($\supseteq$) 假设 $a \in X$。因为 $X \subseteq Y$，所以 $a \in Y$，根据交集的定义，我们有 $a \in X \cap Y$。
\end{itemize}
%% END QUOTATION %%

反过来，假设 $X \cap Y = X$。为了证明 $X \subseteq Y$，令 $a \in X$。因为 $X = X \cap Y$，所以 $a \in X \cap Y$，因此根据交集的定义 $a \in Y$。
\end{cproof}

\begin{exercise}
设 $X$ 为集合。证明 $X \cap \varnothing = \varnothing$。
\end{exercise}

\begin{definition}
\label{defDisjoint}
设 $X$ 和 $Y$ 为集合。如果$X \cap Y$ 为空，则我们说 $X$ 和 $Y$ \textbf{不相交}。
\end{definition}

\begin{example}
集合 $\{ 0,2,4 \}$ 与 $\{ 1,3,5 \}$ 不相交，因为它们没有共同的元素。
\end{example}

\begin{exercise}
设 $a,b,c,d \in \mathbb{R}$ 其中 $a<b$ 且 $c<d$。证明开区间 $(a,b)$ 和 $(c,d)$ 不相交，当且仅当 $b<c$ 或 $d<a$。
\end{exercise}

\subsubsection*{并集 ($\cup$)}
两个集合的\textit{并集}是至少属于其中一个集合的元素的集合。

\begin{definition}
\label{defUnion}
\index{union!pairwise}
设 $X$ 和 $Y$ 为集合。$X$ 和 $Y$ 的（\textbf{成对}）\textbf{并集}，表示为 $X \cup Y$\nindex{union}{$\cup$}{union} \inlatex{cup}\lindexmmc{cup}{$\cup$}, 定义为
\[ X \cup Y = \{ a \mid a \in X \vee a \in Y \} \]
\end{definition}

\begin{example}
设 $E$ 为偶数集合，$O$ 为奇数集合。由于每个整数要么是偶数要么是奇数，$E \cup O = \mathbb{Z}$。注意 $E \cap O = \varnothing$，因此 $\{E,O\}$ 是 $\mathbb{Z}$ \textit{分割} 的一个例子 --- 参见 \Cref{defPartition}。
\end{example}

\begin{exercise}
写出下面集合的元素
\[ \{ 0, 1, 4, 7 \} \cup \{ 1, 2, 3, 4, 5 \} \]
\end{exercise}

\begin{exercise}
将 $[-2,5) \cup [4,7)$ 表示为单个区间。
\end{exercise}

并集运算允许我们定义以下类别的集合，这对我们研究 \Cref{secCountingPrinciples} 中的计数原理特别有用。

\begin{exercise}
设 $X$ 和 $Y$ 为集合。证明 $X \subseteq Y$ 当且仅当 $X \cup Y = Y$。
\end{exercise}

\begin{example}
\label{exIntersectionDistributesOverUnion}
设 $X,Y,Z$ 为集合。证明 $X \cap (Y \cup Z) = (X \cap Y) \cup (X \cap Z)$。
\begin{itemize}
\item ($\subseteq$) 设 $x \in X \cap (Y \cup Z)$。则 $x \in X$，且 $x \in Y$ 或 $x \in Z$。如果 $x \in Y$ 则 $x \in X \cap Y$，如果 $x \in Z$ 则 $x \in X \cap Z$。综合上述两种情况，我们有 $x \in (X \cap Y) \cup (X \cap Z)$。
\item ($\supseteq$) 设 $x \in (X \cap Y) \cup (X \cap Z)$。则 $x \in X \cap Y$ 或 $x \in X \cap Z$。在这两种情况下，根据交集的定义，我们都有 $x \in X$。
第一种情况下我们有 $x \in Y$，第二种情况下我们有 $x \in Z$；综合上述两种情况，我们有 $x \in Y \cup Z$，所以 $x \in X \cap (Y \cup Z)$。
\end{itemize}
\end{example}

\begin{exercise}
\label{exUnionDistributesOverIntersection}
设 $X,Y,Z$ 为集合。证明 $X \cup (Y \cap Z) = (X \cup Y) \cap (X \cup Z)$。
\end{exercise}

\subsubsection*{索引集合族}

我们经常需要取两个以上（甚至无穷多个）集合的交集或并集。例如，我们可能想知道对于每个 $n \ge 1$ 哪些实数是 $[0, 1+\frac{1}{n})$ 的元素，以及哪些实数是其中至少一个集合的元素。

因此，我们现在的任务是将交集和并集的成对概念推广到任意多集合，称为集合的\textit{索引族}。

\begin{definition}
\label{defIndexedFamily}
\index{set!indexed family of}
\index{indexed family}
\index{family of sets}
（\textbf{索引}）\textbf{集合族}是某个\textbf{索引集合} $I$ 中每个元素 $i$ 的集合 $X_i$ 的规范。我们将 $\{ X_i \mid i \in I \}$ 写为索引集合族。
\end{definition}

\begin{example}
\label{exIndexedFamilyOfHalfOpenIntervals}
上面提到的集合 $[0,1+\frac{1}{n})$ 组成一个索引集合族，其索引集合为 $\{ n \in \mathbb{N} \mid n \ge 1 \}$。我们可以将这个集合族缩写为
\[ \{ [0,1+\textstyle\frac{1}{n}) \mid n \ge 1 \} \]
请注意，我们隐含了变量 $n$ 取值范围在自然数上这一事实，并且只是在竖线右侧写下 `$n \ge 1$'，而不是单独定义 $I = \{ n \in \mathbb{N} \mid n \ge 1 \}$ 中并写入 $\{ [0,1+\frac{1}{n}) \mid n \in I \}$。
\end{example}

\begin{definition}
\label{defIndexedIntersectionUnion}
\index{intersection}
\index{union}
\index{intersection!indexed}
\index{union!indexed}
\nindex{intersectionindexed}{$\bigcap_{i \in I}$}{indexed intersection}
\nindex{unionindexed}{$\bigcup_{i \in I}$}{indexed union}
索引族 $\{ X_i \mid i \in I \}$ 的（\textbf{索引}）\textbf{交集}定义为
\[ \bigcap_{i \in I} X_i = \{ a \mid \forall i \in I,\, a \in X_i \} \quad \text{\inlatex{bigcap\_\{i \textbackslash{}in I\}}} \]
索引族 $\{ X_i \mid i \in I \}$ 的（\textbf{索引}）\textbf{并集}定义为
\[ \bigcup_{i \in I} X_i = \{ a \mid \exists i \in I,\, a \in X_i \} \quad \text{\inlatex{bigcup\_\{i \textbackslash{}in I\}}} \]
\end{definition}

\begin{example}
\label{exIndexedIntersectionOfHalfOpenIntervals}
证明半开区间 $[0,1+\frac{1}{n})$ 在 $n \ge 1$ 上的交集是 $[0,1]$。我们将用符号 $\displaystyle \bigcap_{n \ge 1}$ 作为 $\displaystyle \bigcap_{n \in \{ x \in \mathbb{N} ~\mid~ x \ge 1 \}}$ 的简写。

\begin{itemize}
\item $(\subseteq)$ 设 $x \in \displaystyle\bigcap_{n \ge 1} [0,1+\frac{1}{n})$。

则对于所有 $n \ge 1$, $x \in [0,1+\frac{1}{n})$。特别是，$x \ge 0$。

要得到 $x \le 1$，假设$x>1$ --- 我们将得出一个矛盾。由于 $x>1$，我们有 $x-1 > 0$。设 $N \ge 1$ 为大于或等于 $\frac{1}{x-1}$ 的某个自然数，因此 $\frac{1}{N} \le x-1$。那么 $x \ge 1 + \frac{1}{N}$，因此 $x \not\in [0,1+\frac{1}{N})$，与对于所有 $n \ge 1$, $x \in [0,1+\frac{1}{n})$ 的假设矛盾。

所以必有 $x \le 1$，因此 $x \in [0,1]$。

\item ($\supseteq$) 设 $x \in [0,1]$。

为了证明 $x \in \displaystyle \bigcap_{n \ge 1} [0,1+\frac{1}{n})$，我们需要证明对于所有 $n \ge 1$, $x \in [0,1+\frac{1}{n})$ 。所以给定 $n \ge 1$。由于 $x \in [0,1]$，我们有 $x \ge 0$ 和 $x \le 1 < 1+\frac{1}{n}$，因此 $x \in [0,1+ \frac{1}{n})$。
\end{itemize}

因此根据双重包含，$\displaystyle\bigcap_{n \ge 1} [0,1+\frac{1}{n}) = [0,1]$。
\end{example}

\begin{exercise}
将 $\displaystyle\bigcup_{n \ge 1} [0,1-\frac{1}{n}]$ 表示为单个区间。
\end{exercise}

\begin{exercise}
证明 $\displaystyle\bigcap_{n \in \mathbb{N}} [n] = \varnothing$ 和 $\displaystyle\bigcup_{n \in \mathbb{N}} [n] = \{ k \in \mathbb{N} \mid k \ge 1 \}$。
\end{exercise}

索引交集和并集泛化了它们的成对项，正如以下练习所证明的那样。

\begin{exercise}
\label{exIndexedIntersectionUnionGeneralisePairwise}
设 $X_1$ 和 $X_2$ 为集合。证明
\[ X_1 \cap X_2 = \bigcap_{k \in [2]} X_k \quad \text{and} \quad X_1 \cup X_2 = \bigcup_{k \in [2]} X_k \]
\end{exercise}

\begin{exercise}
\label{exSubsetsFiniteIntersectionInhabitedInfiniteIntersectionEmpty}
找到集合 $\{ X_n \mid n \in \mathbb{N} \}$ 的族，使得：
\begin{enumerate}[(i)]
\item $\displaystyle\bigcup_{n \in \mathbb{N}} X_n = \mathbb{N}$；
\item $\displaystyle\bigcap_{n \in \mathbb{N}} X_n = \varnothing$；
\item 对于所有 $i,j \in \mathbb{N}$, $X_i \cap X_j \ne \varnothing$。
\end{enumerate}
\begin{backhint}
\hintref{exSubsetsFiniteIntersectionInhabitedInfiniteIntersectionEmpty}
您需要找到 $\mathbb{N}$ 的一系列子集，使得 (i) 中任何两个子集都有无限多个共同元素，而 (ii) 给定任何自然数，都可以找到其中一个子集\textit{不是}其元素。
\end{backhint}
\end{exercise}

\subsubsection*{补集 ($\setminus$)}

\begin{definition}
\label{defRelativeComplement}
\index{relative complement}
\index{complement!relative}
\nindex{complement}{$\setminus$}{relative complement}
设 $X$ 和 $Y$ 为集合。$Y$ 在 $X$ 上的\textbf{相对补集}，表示为 $X \setminus Y$ \inlatex{setminus}\lindexmmc{setminus}{$\setminus$}，定义为
\[ X \setminus Y = \{ x \in X \mid x \not \in Y \} \]
$\setminus$ 运算也叫\textbf{集合差}运算。有些作者写为 $Y - X$ 而不是 $Y \setminus X$。
\end{definition}

\begin{example}
设 $E$ 为所有偶数的集合。那么 $n \in \mathbb{Z} \setminus E$ 当且仅当 $n$ 是整数且 $n$ 不是偶数；也就是说，当且仅当 $n$ 是奇数。因此 $\mathbb{Z} \setminus E$ 是所有奇数的集合。

此外，$n \in \mathbb{N} \setminus E$ 当且仅当 $n$ 是自然数且 $n$ 不是偶数。由于作为自然数的偶整数正是偶自然数，因此 $\mathbb{N} \setminus E$ 正是所有奇自然数的集合。
\end{example}

\begin{exercise}
写出下面集合的元素
\[ \{ 0, 1, 4, 7 \} \setminus \{ 1, 2, 3, 4, 5 \} \]
\end{exercise}

\begin{exercise}
将 $[-2,5) \setminus [4,7)$ 和 $[4,7) \setminus [-2,5)$ 表示为单个区间。
\end{exercise}

\begin{exercise}
\label{exSetMinusSetMinus}
设 $X$ 和 $Y$ 为集合。证明 $Y \setminus (Y \setminus X)= X \cap Y$，并推导出 $X \subseteq Y$ 当且仅当 $Y \setminus (Y \setminus X) = X$。
\end{exercise}

\subsubsection*{与逻辑运算符和量词的比较}

精明的读者会注意到集合运算与我们在 \Cref{chLogicalStructure} 中看到的逻辑运算符和量词之间的某些相似之处。

事实上，这可以总结为下表。每一行中，两列表达式是等效的，其中 $p$ 表示 `$a \in X$'，$q$ 表示`$a \in Y$'，$r(i)$ 表示`$a \in X_i$'。

\begin{center}\begin{tabular}{c|c}
集合 & 逻辑 \\ \hline
$a \not\in X$ & $\neg p$ \\
$a \in X \cap Y$ & $p \wedge q$ \\
$a \in X \cup Y$ & $p \vee q$ \\
$a \in \bigcap_{i \in I} X_i$ & $\forall i \in I,\, r(i)$ \\
$a \in \bigcup_{i \in I} X_i$ & $\exists i \in I,\, r(i)$ \\
$a \in X \setminus Y$ & $p \wedge (\neg q)$
\end{tabular}\end{center}

逻辑和集合论之间的转换不止于此。事实上，正如下面定理所示，逻辑运算符 (\Cref{thmDeMorganLogicalOperators}) 和量词 (\Cref{thmDeMorganQuantifiers}) 的德摩根定律也适用于并集和交集的集合运算。

\begin{theorem}[集合的德摩根定律]
\label{thmDeMorganForSets}
\index{de Morgan's laws!for sets}
给定集合 $A,X,Y$ 和族 $\{ X_i \mid i \in I \}$，我们有
\begin{enumerate}[(a)] 
\item $A \setminus (X \cup Y) = (A \setminus X) \cap (A \setminus Y)$;
\item $A \setminus (X \cap Y) = (A \setminus X) \cup (A \setminus Y)$;
\item $\displaystyle A \setminus \bigcup_{i \in I} X_i = \bigcap_{i \in I} (A \setminus X_i)$;
\item $\displaystyle A \setminus \bigcap_{i \in I} X_i = \bigcup_{i \in I} (A \setminus X_i)$.
\end{enumerate}
\end{theorem}

\begin{cproof}[of (a)]
设 $a$ 为任意的元素。根据并集和补集的定义，$a \in A \setminus (X \cup Y)$ 等价于逻辑公式
\[ a \in A \wedge \neg (a \in X \vee a \in Y) \]
根据逻辑运算符的德摩根定律，这等价于
\[ a \in A \wedge (a \not\in X \wedge a \not\in Y) \]
反过来，这等价于
\[ (a \in A \wedge a \not\in X) \wedge (a \in A \wedge a \not\in Y) \]
但根据交集和补集的定义，这相当于
\[ a \in (A \setminus X) \cap (A \setminus Y) \]
因此 $A \setminus (X \cup Y) = (A \setminus X) \cap (A \setminus Y)$。
\end{cproof}

\begin{exercise}
完成集合的德摩根定律的证明。
\end{exercise}

\subsubsection*{乘积 ($\times$)}

\begin{definition}
\label{defCartesianProductPairwise}
\index{product of sets!pairwise}
\index{ordered pair}
设 $X$ 和 $Y$ 为集合。 $X$ 和 $Y$ 的（\textbf{成对}）\textbf{笛卡尔积} 是集合 $X \times Y$\nindex{cartesian product}{$\times$}{cartesian product} \inlatex{times}\lindexmmc{times}{$\times$} 定义为
\[ X \times Y = \{ (a,b) \mid a \in X \wedge b \in Y \} \]
元素 $(a,b) \in X \times Y$ 称为 \textbf{有序对}，其定义属性为，对于所有 $a,x \in X$ 和所有 $b,y \in Y$ ，当且仅当 $a=x$ 且 $b=y$ 时，我们有 $(a,b) = (x,y)$。
\end{definition}

\begin{example}
如果您学过微积分，您可能对集合 $\mathbb{R} \times \mathbb{R}$ 很熟悉。
\[ \mathbb{R} \times \mathbb{R} = \{ (x,y) \mid x,y \in \mathbb{R} \} \]
形式上，这是一组有序的实数对。从几何角度来说，如果我们将 $\mathbb{R}$ 解释为无限延伸的直线，则集合 $\mathbb{R} \times \mathbb{R}$ 是（实）平面：元素 $(x,y) \in \mathbb{R} \times \mathbb{R}$ 描述了平面上坐标为 $(x,y)$ 的点。

我们可以更进一步。例如，下面的集合：
\[ \mathbb{R} \times \{ 0 \} = \{ (x,0) \mid x \in \mathbb{R} \} \]
正是 $x$ 轴。我们可以将图描述为 $\mathbb{R} \times \mathbb{R}$ 的子集。事实上，$y=x^2$ 的图由下面式子给出
\[ G = \{ (x,y) \in \mathbb{R} \times \mathbb{R} \mid y = x^2 \} = \{ (x,x^2) \mid x \in \mathbb{R} \} \subseteq \mathbb{R} \times \mathbb{R} \]
\end{example}

\begin{exercise}
写出集合 $\{ 1, 2 \} \times \{ 3, 4, 5 \}$ 的所有元素。
\end{exercise}

\begin{exercise}
设 $X$ 为集合。证明 $X \times \varnothing = \varnothing$。
\end{exercise}

\begin{exercise}
\label{exCartesianProductNotAssociative}
设 $X$, $Y$ 和 $Z$ 为集合。在什么条件下 $X \times Y = Y \times X$ 成立？在什么条件下 $(X \times Y) \times Z = X \times (Y \times Z)$ 成立？
\end{exercise}

我们可能需要计算两个以上集合的笛卡尔积。例如，无论集合 $\mathbb{R} \times \mathbb{R} \times \mathbb{R}$ 是什么，它的元素\textit{应该}是由元素 $a,b,c \in \mathbb{R}$ 构成的有序三元组 $(a,b,c)$。这就是下面定义派上用场的地方。

\begin{definition}
\label{defCartesianProductNFold}
\index{product of sets!nfold@$n$-fold}
\index{ordered $n$-tuple}
设 $n \in \mathbb{N}$ 并设 $X_1, X_2, \dots, X_n$ 为集合。$X_1, X_2, \dots, X_n$ 的（\textbf{$n$ 折}）\textbf{笛卡尔积}是集合 $\prod_{k=1}^n X_k$\nindex{cartesian product nfold}{$\Pi_{k=1}^n$}{cartesian product ($n$-fold)} \inlatex{prod\_\{k=1\}\^{}\{n\}}\lindexmmc{prod}{$\Pi_{k=1}^{n}$} 定义为 
\[ \prod_{k=1}^n X_k = \{ (a_1, a_2, \dots, a_n) \mid a_k \in X_k \text{ for all } 1 \le k \le n \} \]
元素 $(a_1, a_2, \dots, a_n) \in \prod_{k=1}^n X_k$ 称为\textbf{有序 $k$ 元组}，其定义属性为，对于所有 $1 \le k \le n $ 和所有 $a_k,b_k \in X_k$，我们有 $(a_1, a_2, \dots, a_n) = (b_1, b_2, \dots, b_n)$ 当且仅当对于所有 $1 \ le k \le n$, $a_k = b_k$。

给定一个集合$X$，$X^n$\nindex{cartesian product nfold}{$X^n$}{cartesian product ($n$-fold)} 表示集合 $\prod_{k=1}^n X$。有时我们也会写做
\[ X_1 \times X_2 \times \cdots \times X_n = \prod_{k=1}^n X_k \]
\end{definition}

\begin{example}
在 \Cref{exCartesianProductNotAssociative} 中您可能已经注意到集合 $(X \times Y) \times Z$ 和 $X \times (Y \times Z)$ 并不总是相等 --- \Cref{defCartesianProductNFold} 引入了第三个 $X$、$Y$ 和 $Z$ 可能不相等的笛卡尔积。例如，考虑当 $X=Y=Z=\mathbb{R}$ 时。则
\begin{itemize}
\item $(\mathbb{R} \times \mathbb{R}) \times \mathbb{R}$ 的元素是有序对 $((a,b),c)$，其中 $(a,b)$ 本身是一对有序实数，$c$ 是一个实数。
\item $\mathbb{R} \times (\mathbb{R} \times \mathbb{R})$ 的元素是有序对 $(a,(b,c))$，其中 $a$ 是实数， $(b,c)$ 是一对有序实数。
\item $\mathbb{R} \times \mathbb{R} \times \mathbb{R}$ ($=\mathbb{R}^3$) 的元素是有序三元组 $(a,b,c)$，其中 $a$、$b$ 和 $c$ 是实数。
\end{itemize}
因此，尽管这三个集合\textit{看起来}相同，但仔细观察定义就会发现它们之间存在细微差异。从某种意义上说，它们的相同之处在于它们之间存在\textit{双射} --- 双射的概念将在\Cref{secInjectionsSurjections}中引入。
\end{example}

\begin{tldr}{secSetOperations}

\subsubsection*{二元集合运算}

\begin{tldrlist}
\tldritem{defIntersection} 集合 $X$ 和 $Y$ 的\textit{交集}为 $X \cap Y = \{ a \mid a \in X \wedge a \in Y \}$。
\tldritem{defUnion} 集合 $X$ 和 $Y$ 的\textit{并集}为 $X \cup Y = \{ a \mid a \in X \vee a \in Y \}$。
\tldritem{defRelativeComplement} 集合 $Y$ 在集合 $X$ 上的\textit{相对补集}为 $X \setminus Y = \{ a \mid a \in X \wedge a \not\in Y \}$。
\tldritem{defCartesianProductPairwise} 集合 $X$ 和 $Y$ 的\textit{笛卡尔积}为 $X \times Y = \{ (a,b) \mid a \in X \wedge b \in Y \}$，其中 $(a,b)$ 是一个\textit{有序对}。
\end{tldrlist}

\subsubsection*{索引集合运算}

\begin{tldrlist}
\tldritem{defIndexedFamily} \textit{索引集合族} $\{ X_i \mid i \in I \}$ 是某个\textit{索引集合} $I$ 中每个元素 $i$ 的集合 $X_i$ 的规范。
\tldritem{defIndexedIntersectionUnion} 族 $\{ X_i \mid i \in I \}$ 的\textit{索引交集}为 $\bigcap_{i \in I} X_i = \{ a \mid \forall i \in I,\, a \in X_i \}$；该族的 \textit{索引并集} 为 $\bigcup_{i \in I} X_i = \{ a \mid \exists i \in I,\, a \in X_i \}$。
\tldritem{defCartesianProductNFold} $[n]$ 索引族 $\{ X_i \mid i \in [n] \}$ 的 \textit{$n$ 次笛卡尔积}是集合   $\prod_{k=1}^n X_k = \{ (a_1, \dots, a_n) \mid a_k \in X_k \textbf{ for all } k \in [n] \}$，其中 $(a_1, \dots, a_n)$ 为 \textit{有序 $n$ 元组}。
\end{tldrlist}

\end{tldr}

\index{set|)}